% TOP_PROBLEM: {"id": 3128, "title": "Switching reconstruction conjecture", "category_id": 3, "category": {"id": 3, "name": "graph_theory", "display_name": "Graph Theory", "description": "Problems involving graphs, networks, and their properties.", "slug": "graph-theory", "order_index": 3, "created_at": "2026-07-31T15:26:25.671Z"}, "status": "open", "problem_number": "OPG-46951", "set_id": 12}
% FIRST_POSTED: 2026-10-01
% ENTRY_KIND: statement_only
% UnsolvedMath ID 3128; source catalogue code OPG-46951
% !TeX program = lualatex
% Definition and sources only; this file is not an LLM solution attempt.
\documentclass[11pt,a4paper]{article}
\usepackage[margin=25mm]{geometry}
\usepackage{fontspec}
\setmainfont{lmroman10-regular.otf}[BoldFont=lmroman10-bold.otf,
 ItalicFont=lmroman10-italic.otf,BoldItalicFont=lmroman10-bolditalic.otf]
\usepackage{amsmath,amssymb,mathtools}
\usepackage{unicode-math}
\setmathfont{latinmodern-math.otf}
\AtBeginDocument{\let\setminus\smallsetminus}
\usepackage{xurl,hyperref}
\hypersetup{colorlinks=true,urlcolor=blue}
\setcounter{secnumdepth}{0}
\setlength{\parindent}{0pt}
\setlength{\parskip}{5pt}
\setlength{\emergencystretch}{3em}
\begin{document}
\section*{Switching reconstruction conjecture}\label{3128}
{\small UnsolvedMath ID 3128 · OPG-46951 · Graph Theory · OpenGarden}
\subsection{Definitions and mathematical statement}
Let $G=(V,E)$ be a finite simple undirected graph. For a vertex $v\in V$,
its vertex switching $G^{(v)}$ retains the same vertex set, leaves all
adjacencies within $V\setminus\{v\}$ unchanged, and replaces the neighbour
set of $v$ by its complement within $V\setminus\{v\}$. Thus every edge
incident with $v$ becomes a nonedge and every incident nonedge becomes
an edge; no loop is introduced.

With $[H]$ denoting an unlabelled isomorphism class, define the switching
deck to be the multiset
\[
 \mathcal S(G)=\{\!\{[G^{(v)}]:v\in V\}\!\}.
\]
There is one card for each single-vertex switching, and equal isomorphism
classes are kept with multiplicity. The switched vertex is not marked
in its card, nor is there a common labelling supplied across cards.

The conjecture asserts that for any finite simple graphs $G,H$ on
at least five vertices,
\[
 \mathcal S(G)=\mathcal S(H)\quad\Longrightarrow\quad G\cong H.
\]
A card has the same number of vertices as its original graph; this
operation is different from vertex deletion. Neither connectedness nor
any degree condition is assumed.

The deck contains the results of switching once at each individual
vertex, not all graphs obtainable by arbitrary sequences of switchings.
The target is ordinary graph isomorphism, not merely membership in the
same switching-equivalence class. Consequently both the single-switch
convention and card multiplicities are necessary to fix the exact
reconstruction question.
\subsection{Short English statement}
Is every simple graph on at least five vertices determined by the multiset of graphs obtained by complementing all adjacencies at one vertex?
\subsection{Sources}
\begin{itemize}
\item[S1] ULAM.AI and the UnsolvedMath Contributors, supplied problems.json, record 3128 (OPG-46951), OpenGarden collection. Catalogue identity and source transcription; CC BY 4.0.
\url{https://huggingface.co/datasets/ulamai/UnsolvedMath}.
\item[S2] Open Problem Garden, Switching reconstruction conjecture. Original problem and discussion; the mathematical formulation below is expanded from the supplied source record.
\url{http://www.openproblemgarden.org/op/switching_reconstruction_conjecture}.
\end{itemize}
\end{document}
